\paragraph{Section Overview}
This section summarizes what the current Neural Computer (NC) prototypes can already do and where they still fail as general-purpose systems.
We then define Completely Neural Computers (CNCs) using operational requirements and map those requirements to measurable challenge families.
Finally, we position NCs relative to world models and computer-use agents, and we outline design hypotheses beyond the current video-based instantiation.

\subsection{From Neural Computers to Completely Neural Computers}

\paragraph{Current Status of NCs}
Our CLI and GUI prototypes show that a learned latent runtime state can reproduce core interface behaviors with measurable fidelity.
In terminal environments, character-level OCR accuracy reaches $0.54$ (\Cref{tab:cligen-ocr}), while desktop cursor control reaches $98.7\%$ with explicit visual supervision (\Cref{tab:cursor-loss}).
In GUIWorld, $110$ hours of goal-directed trajectories outperform approximately $1{,}400$ hours of random exploration (\Cref{tab:gui-data-quality}), underscoring the role of aligned supervision.
These results support a bounded claim: current NCs can learn interface I/O and short-horizon control at moderate training cost when data are well aligned.
They do not yet establish robust long-horizon reasoning, stable routine reuse, or deployment-grade governance.

To move from prototype NCs to general-purpose systems, four properties are central.
An NC instance should be computationally universal~\citep{turing1936computable,siegelmann1992computational}, universally programmable~\citep{von1993first,wilkes1981best}, behavior-consistent unless explicitly updated~\citep{queloz2025explainability}, and able to realize practical advantages of neural architectural and programming-language semantics.
While several sequential neural architectures are universal in principle~\citep{siegelmann1992computational,perez2021attention}, reliably programming trained instances remains difficult.
Preliminary attempts such as Neural Virtual Machine and NeuroLISP suggest feasibility but also highlight engineering gaps~\citep{katz2019programmable,davis2022neurolisp}.
Long-horizon stability remains an open systems problem in neural runtimes~\citep{kirkpatrick2017overcoming,calanzonelogically}.
Section~\ref{sec:roadmap} reframes these requirements as operational targets.
Existing world-model work rarely frames these targets in computability and programmability terms, and we revisit this relation in Section~\ref{sec:relations}.

\paragraph{Fundamental differences between NCs and conventional computers}
We compare NCs with conventional computers.
Here, \emph{conventional computers} denote random-access machines with instruction-set architecture and layered OS/application stacks programmed through human-designed high-level languages~\citep{hartmanis1974power,anagnostopoulos1973computer,backus1957fortran,backus1963revised}.
NCs differ from this paradigm in both architectural semantics and programming-language semantics.

At the architectural level, conventional machines implement local, compositional symbolic semantics, which provide exactness and interpretability but can be brittle under noise and model mismatch~\citep{Newell1976,BROOKS1991139}.
NCs use holistic, distributed numerical semantics, trading strict local exactness for robustness and generalization~\citep{Hinton1986,Bishop2006,lecun2015deep}.
Such representations are particularly effective in domains with high-dimensional observations, soft constraints, and globally coupled dynamics~\citep{bengio2013representation,Smolensky1988-SMOOTP-2,silver2017mastering,vaswani2017attention}.
Although conventional systems can emulate NC behavior in principle, that emulation can introduce avoidable conceptual and engineering overhead when the task is naturally neural.

At the programming-language level, NC semantics are learned from data rather than fully specified by human-designed grammar and runtime rules.
Prompts and interaction traces can therefore function as executable specifications~\citep{Reynolds2021}.
This makes natural-language and multimodal intent a practical programming interface at scale, which has historically been hard to realize in non-neural systems~\citep{jm3,wei2023chainofthoughtpromptingelicitsreasoning}.

\paragraph{Definition of Completely Neural Computers}
We define a Neural Computer instance as \emph{complete} if it is (1) Turing complete, (2) universally programmable, (3) behavior-consistent unless explicitly reprogrammed, and (4) able to realize practically useful neural architectural and programming-language semantics relative to conventional systems.
The next section turns these properties into concrete challenge families and evaluation signals.

\subsection{A Roadmap Towards CNC}
\label{sec:roadmap}
We organize the CNC roadmap into five challenge families.
The first three correspond directly to universality, programmability, and behavior consistency.
The last two unpack the fourth property into architectural semantics and programming-language semantics.

\paragraph{Turing completeness}
In the limit, a CNC should be as expressive as a universal computing model~\citep{turing1936computable,siegelmann1992computational}.
Any physical implementation is finite, so we treat universality as an asymptotic design requirement rather than a near-term pass/fail benchmark.
Operationally, the key question is scaling behavior: as effective memory and context grow, can one NC instance execute longer algorithmic procedures without qualitatively new failure modes?
Prior work on recurrent neural networks (RNNs), Neural Turing Machines (NTMs), and Differentiable Neural Computers (DNCs) motivates this direction, while also exposing practical limits under finite precision and finite memory~\citep{graves2014neural,graves2016hybrid,perez2021attention,rusu2016progressive,vaswani2017attention}.

\paragraph{Universal programmability}
Universal programmability asks whether input sequences can install reusable routines instead of only triggering one-off behaviors.
The practical target is routine lifecycle control: install, invoke, compose, and retain routines under distribution shift.
Useful proxies include installation success, cross-task reuse without retraining, and compositional generalization of previously installed routines.
Compositional neural routines are a plausible path toward this objective~\citep{reed2015neural}.

\paragraph{Behavior consistency}
A CNC should preserve function during ordinary execution and change function only through explicit updates~\citep{queloz2025explainability}.
This requires a runtime contract that separates function use from function update.
Operationally, every behavioral change should be attributable to a versioned update and reversible through replay/rollback.
Progress should be tracked by same-version reproducibility, update traceability, rollback success, and drift rates in long-horizon workloads.

\paragraph{Architectural semantics}
Because NC behavior is parameterized in continuous space, learning can produce mappings beyond observed training pairs and can introduce new functional primitives~\citep{poggio2019theoretical,ha2016hypernetworks}.
Recombining these primitives can yield out-of-distribution functional generalization~\citep{lake2018generalization}.
For CNCs, the critical test is practical transfer rather than short demos: after learning many trajectories of an interface family, can the system generalize to unseen tasks, files, and contexts with stable behavior?
Beyond emulating symbolic APIs, NCs may natively support uncertainty-aware inference, dense retrieval, and tightly coupled perception-control loops that operate directly on distributed states~\citep{marcus2018deep,kingma2013auto,bengio2013representation,graves2016hybrid,silver2017mastering,clements2019estimating,parisi2019continual}.

Because NC memory is numerical, computer configuration can be cast as optimization over internal state rather than only authoring discrete symbolic code.
Depending on objective differentiability, both gradient-based and gradient-free methods are relevant~\citep{adam2014method,wierstra2014natural}.
This reframes part of system construction as synthesizing target behaviors through state optimization, with progress measured by convergence reliability and stability relative to combinatorial program search~\citep{innes2019differentiable,hong2023metagpt}.

\paragraph{Programming-language semantics}
Learned programming-language semantics shift development from rigid code authoring to specification by instructions, demonstrations, and constraints~\citep{brown2020language,wei2022chain,ouyang2022training}.
Brief inputs can then trigger longer structured behaviors when routines are already installed.
Development emphasis moves toward curating data, specifying intent, and verifying behavior under learned semantics rather than hard-coding every execution path~\citep{austin2021program}.

This shift is supported by data availability.
Paired I/O traces are continuously generated during ordinary computer use and are often far more abundant than high-quality hand-written programs~\citep{flener2008introduction,yao2022webshop}.
These traces act as executable specifications of user intent and environment response~\citep{cypher1993watch,yao2022react}.
This asymmetry in supervision favors NC training regimes that exploit ubiquitous interaction logs while preserving rigorous evaluation.

\subsection{Relations to Other Foundation Models}
\label{sec:relations}
\paragraph{World Models}
World models learn environment dynamics through action-conditioned prediction~\citep{ha2018world}.
Their target environments range from broad real-world sensory streams to narrow control systems~\citep{schmidhuber1990making,feng2023finetuning}.
From this perspective, NCs are world models specialized for interactive computers.
The key distinction is objective: predictive coherence is necessary, but CNCs additionally require installable routines, reusable behavior, and governed updates.
Many current world-model pipelines also rely heavily on computer-generated data, which creates methodological overlap with NC development~\citep{richter2016playing,tobin2017domain,dosovitskiy2017carla,garcia2023robust}.

\paragraph{Computer Use Agents}
Computer-use agents (CUAs) combine foundation models with conventional computers through an external control layer~\citep{openaiComputerUsingAgent,claudeComputerTool,sager2025comprehensivesurveyagentscomputer}.
This design leverages both model-level cognition and the precision of conventional software stacks~\citep{song2025coact1computerusingagentscoding}.
Its main constraints come from low-bandwidth I/O mediation and strict separation of internal states between model and runtime.
As a result, application development often remains anchored to fixed symbolic semantics on the conventional side, which can limit end-to-end optimization over richer latent representations~\citep{liang-etal-2017-neural,hu2017learningreasonendtoendmodule,gaunt2016terpretprobabilisticprogramminglanguage}.
Near term, CUAs and NCs should be viewed as complementary: CUAs remain strong production baselines, while NCs investigate what can be internalized into a single learned runtime.

\subsection{Additional Thoughts}

\paragraph{Video models as a pragmatic prototype substrate}
We currently build NC prototypes on top of state-of-the-art video models because they provide the simplest end-to-end substrate for joint pixel dynamics and action-conditioned control.
This choice is pragmatic, not fundamental.
In our experiments, symbolic and algorithmic reasoning in terminal settings remains inconsistent for most strong video models, and even simple arithmetic can fail (\Cref{tab:cligen-arith}).
One frontier model in our probe reaches $71\%$ on arithmetic, indicating partial progress but not deployment-grade reliability.
At the same time, we do not claim that video models cannot reason more broadly; recent work reports zero-shot reasoning in naturalistic settings~\citep{wiedemer2025video}.
Our current evidence suggests that CNC-level reliability will likely require additional architectural and training ingredients beyond scaling today's generic video generators.

\paragraph{A hypothesis: machine-native neural architectures}
This subsection states a conjecture rather than an experimental conclusion.
Closing the reasoning gap may require machine-native neural designs instead of architectures primarily inherited from biological perception analogies.
Many influential architectures are highly engineered yet still rely on continuous distributed representations in which reasoning emerges implicitly~\citep{fukushima1980neocognitron,lecun2015deep,vaswani2017attention}.
CNCs may benefit from designs that explicitly incorporate discrete operations, compositional structure, and verifiable computation while remaining compatible with neural optimization.

\paragraph{Neural network generation via NC interaction}
Generating neural networks can be viewed as a programming activity that synthesizes both architecture and weights.
Because NC semantics are already neural and numerical, such generation can operate directly over internal memory states rather than through symbolic translation layers.
NCs can also be programmed through interaction: sequences of inputs, observations, and outcomes act as executable specifications that shape internal routines~\citep{cypher1993watch}.
This suggests a path where users iteratively generate and refine neural modules through interaction traces.

\paragraph{Unified hardware requirements and data representation}
NCs treat tensors and tensor-to-tensor transforms as primary computational primitives.
This contrasts with conventional stacks that combine many heterogeneous abstractions (files, sockets, processes, pointers, and subsystem-specific APIs) coordinated by OS services.
A tensor-uniform pipeline can simplify optimization and tooling because compilers, profilers, and runtime systems target a shared representation.
The same representation can naturally support multimodal computation across vision, language, audio, control, and planning.
